% !TeX root = ../cosmology_summary.tex

%% --------------------------------------------------------
\chapter{Програмні пакети Python для космологічних досліджень}
\label{sec:python_packages}
%% --------------------------------------------------------

Сучасна теоретична та спостережна космологія є високотехнологічною дисципліною, де перевірка теоретичних моделей та аналіз даних експериментів (таких як \emph{Planck}, \emph{Euclid} чи \emph{JWST}) спираються на потужний інструментарій чисельного моделювання. Найбільшого поширення в академічному середовищі набула екосистема мови Python, яка об'єднує спеціалізовані космологічні пакети та загальнонаукові бібліотеки для аналізу даних і машинного навчання.

Нижче наведено та класифіковано основні програмні пакети, які є стандартом у сучасних космологічних дослідженнях:

\begin{enumerate}[leftmargin=*]
    \item \textbf{Пакети для лінійного аналізу збурень та еволюції СМВ:}
    \begin{itemize}
        \item \href{http://class-code.net/}{\texttt{CLASS}} (\emph{Cosmic Linear Anisotropy Solving System}) --- один із двох найпопулярніших кодів для інтегрування лінійних збурень у ранньому Всесвіті. Написаний на мові C, але має повну та зручну обгортку \texttt{classy} для Python. Пакет дозволяє з високою точністю розраховувати спектри потужності температурних анізотропій та поляризації реліктового випромінювання, а також спектри потужності матерії $P(k)$.
        \item \href{https://camb.readthedocs.io/}{\texttt{CAMB}} (\emph{Code for Anisotropy in the Microwave Background}) --- класичний інструмент для розрахунку анізотропії CMB. Його обчислювальне ядро реалізоване на Fortran, проте Python-версія повністю інтегрована в сучасні конвеєри обробки даних і дозволяє гнучко модифікувати рівняння стану темної енергії або параметри первинного спектра збурень.
    \end{itemize}

    \item \textbf{Статистичний аналіз та оцінка параметрів (MCMC sampling):}
    \begin{itemize}
        \item \href{https://cobaya.readthedocs.io/}{\texttt{Cobaya}} --- сучасний фреймворк для байєсівського аналізу космологічних даних, розроблений для роботи з архівами \emph{Planck} та майбутніх місій. \texttt{Cobaya} дозволяє зв'язувати розрахунки з \texttt{CAMB} або \texttt{CLASS} та запускати алгоритми Монте-Карло за схемами марковських ланцюгів (MCMC) для пошуку обмежень на космологічні параметри ($\Omega_m$, $H_0$, $n_s$ тощо).
        \item \href{https://github.com/brinckmann/montepython_public}{\texttt{MontePython}} --- популярний MCMC-семплер, який традиційно використовується у зв'язці з кодом \texttt{CLASS} для порівняння теоретичних моделей із космологічними спостереженнями та побудови контурів довірчих інтервалів.
    \end{itemize}

    \item \textbf{Загальні астрофізичні фреймворки та обробка даних:}
    \begin{itemize}
        \item \href{https://www.astropy.org/}{\texttt{Astropy}} --- фундаментальна бібліотека для астрономічних досліджень. Модуль \texttt{astropy.cosmology} містить готові класи для швидкого розрахунку різних типів відстаней (супутньої, кутової, світності) та еволюції густини енергії для стандартних моделей ($\Lambda\text{CDM}$, $w\text{CDM}$).
        \item \href{https://healpy.readthedocs.io/}{\texttt{healpy}} --- Python-обгортка для бібліотеки HEALPix (\emph{Hierarchical Equal Area isoLatitude Pixelization}). Це критично важливий інструмент для роботи з картами реліктового випромінювання та великомасштабної структури, розбитими на сфері. Пакет дозволяє виконувати сферичне гармонічне перетворення для отримання мультипольних коефіцієнтів $C_\ell$.
    \end{itemize}
\end{enumerate}

Для чисельного розв'язання нестандартних динамічних систем (наприклад, для точного інтегрування рівняння Кляйна-Ґордона-Фока зі складними видами потенціалів інфлатона на етапі розігріву) зазвичай використовуються базові наукові пакети \href{https://numpy.org/}{\texttt{NumPy}} та \href{https://scipy.org/}{\texttt{SciPy}} (зокрема, модулі для розв'язання жорстких систем звичайних диференціальних рівнянь \texttt{scipy.integrate}). Візуалізація отриманих спектрів та еволюційних треків традиційно виконується за допомогою бібліотеки \href{https://matplotlib.org/}{\texttt{Matplotlib}}.